\begin{abstract}
Games, film, robot learning and embodied AI all depend on 3D
environments that can be simulated, in which every object is a body of
its own with collision, mass and joints, stands under the gravity of the
ground beneath it, can be changed afterwards, and belongs to the person
who made it. Current generative methods rarely deliver this: world
models produce frames that no simulator can step, and one-shot 3D
generators produce a single fused scene without separate bodies or clear
rights to the result. We present SCORE, a framework that synthesises
such environments by building them the way film studios build scenes in
OpenUSD. Each scene sits within its parent scene, and each kind of work
on it (layout, surfaces, physics, light, motion, effects, sound and
people) is a layer of its own that is checked and changed independently
of the others. The creator locks the design from the top down, first an
environment concept that fixes its style and its physics and then a
concept for each scene; open vision foundation models, procedural tools
and coding agents then make every layer below, each checked by measured
gates before the next is made. Every model in the framework allows
commercial use of its output, so the creator remains the author, and
every step is keyed on its inputs, so an edit reruns only the steps that
depend on it. The result is a single OpenUSD environment whose physics a
simulator reads as written; Newton, with MuJoCo's solver, steps it, and
each character in it receives a sim body on Newton driven by its SOMA
segment rotations. We describe the framework and its method, with a
first environment as the evidence; its four demonstrations are marked as
gaps until its batch has run. We close with the route toward simulated
bodies that are true to the individual person and trained models that
drive them. Code and paper are available at
https://github.com/Babon-Innovations-b-v/score (MIT licence).
\end{abstract}
